\documentclass[11pt]{article}

\usepackage{amsmath}
\usepackage{hyperref}

\title{Zanistarast: A Falsifiable Ontology-Based Civilizational Model}
\author{Veysi yê MALA SAF}
\date{}

\begin{document}

\maketitle

\begin{abstract}
This paper introduces a falsifiable, ontology-based civilizational framework integrating epistemology, economic structure, and ethical coherence into a unified system. The model is designed to be machine-readable, logically structured, and empirically testable.
\end{abstract}

\section{Introduction}

Modern systems are fragmented between knowledge, structure, and reality.

This model proposes a unified framework based on structural coherence.

\section{Framework}

Layers:

\begin{itemize}
\item Ontology (Hebûn)
\item Epistemology (Zanabûn)
\item Economy (Mabûn)
\item Method (Rasterast)
\item Synthesis (Zanistarast)
\end{itemize}

\section{Core Model}

\[
Stability = f(Alignment, Coherence, Responsibility, Entropy^{-1})
\]

\section{Logic}

\begin{itemize}
\item Ontology precedes knowledge
\item Knowledge requires validation
\item Economy requires responsibility
\item Claims require filtering
\end{itemize}

\section{Falsifiability}

The model fails if:

\begin{itemize}
\item Alignment does not affect stability
\item Responsibility does not reduce entropy
\item Coherence does not increase durability
\end{itemize}

\section{Conclusion}

This framework represents a transition from narrative-based systems to structural civilizational engineering.

\end{document}


